\documentclass[11pt]{article}

\usepackage[a4paper,top=2.35cm,bottom=2.45cm,left=2.15cm,right=2.15cm]{geometry}
\usepackage{fontspec}
\setmainfont{TeX Gyre Termes}
\usepackage{amsmath,amssymb}
\usepackage{booktabs,tabularx,array}
\usepackage{graphicx}
\usepackage{xcolor}
\usepackage[font=small,labelfont=bf]{caption}
\usepackage[section]{placeins}
\usepackage[round,authoryear]{natbib}
\usepackage{microtype}
\usepackage{xspace}
\usepackage{lineno}
\usepackage{hyperref}
\hypersetup{
  colorlinks=true,
  linkcolor=black,
  citecolor=blue!45!black,
  urlcolor=blue!55!black,
  pdftitle={Background suppression by a chimney-mounted TES focal plane and BGO anticoincidence},
  pdfauthor={Authors omitted for review}
}

\renewcommand{\arraystretch}{1.13}
\setlength{\emergencystretch}{3em}
\setlength{\parindent}{1.25em}
\setlength{\parskip}{0pt}
\raggedbottom
\newcolumntype{Y}{>{\raggedright\arraybackslash}X}
\newcommand{\keV}{\ensuremath{\,\mathrm{keV}}\xspace}
\newcommand{\MeV}{\ensuremath{\,\mathrm{MeV}}\xspace}
\newcommand{\cm}{\ensuremath{\,\mathrm{cm}}\xspace}
\newcommand{\cps}{\ensuremath{\,\mathrm{s^{-1}}}\xspace}
\newcommand{\phflux}{\ensuremath{\,\mathrm{ph\,cm^{-2}\,s^{-1}}}\xspace}
\newcommand{\aeff}{\ensuremath{A_{\mathrm{eff}}}}

\title{\textbf{Background suppression by a chimney-mounted TES focal plane and BGO anticoincidence for a balloon-borne 511 keV Laue telescope}}
\author{Authors omitted for review\\[0.25em]\small Affiliations omitted for review}
\date{}

\begin{document}
\maketitle
\linenumbers

\begin{abstract}
A balloon-borne Laue telescope can concentrate 511\keV photons onto a compact
transition-edge-sensor focal plane, but detector and cryostat backgrounds can
erase the benefit of focusing. We compare a reference detector assembly with a
geometry that translates the transition-edge-sensor array laterally into an
aluminum side chimney and surrounds it with a three-part bismuth germanate
anticoincidence shield. Both assemblies are evaluated with the same
eight-family atmospheric radiation field, activation treatment, detector
response, event selection, and 20-day synthetic trajectory. The matched
simulation finds that the fraction of selected delayed events originating in
the central cold assembly falls from 32.18\% to 3.92\%. In the
chimney geometry, a 50\keV offline bismuth-germanate veto reduces the prompt and
delayed rates in the 510.58--511.42\keV window from 2.235 to 0.00579\cps and
from 0.1275 to 0.00402\cps, respectively. The final focused-signal effective
area is 15.085$\pm$0.045\,cm$^2$, 99.75\% of the reference value. At the
day-15 trajectory state, the total background falls from 0.05128 to
0.009280\cps. Over 20 days, a top-of-atmosphere line flux of
$10^{-4}\phflux$ gives 1678 signal
counts over 15552 background counts and an Asimov significance of 13.22. The
Gaussian 3$\sigma$ line-flux threshold is
$(2.229\pm0.209)\times10^{-5}\phflux$, a factor of 2.43 below the reference.
The matched result is consistent with geometric decoupling of the focal plane
from activation-prone cold structures while retaining BGO coverage and focused
throughput.
\end{abstract}

\noindent\textbf{Keywords:} 511 keV annihilation line; balloon-borne telescope;
Laue lens; transition-edge sensor; bismuth germanate; activation background

\section{Introduction}
\label{sec:introduction}

The 511\keV line from electron--positron annihilation is a long-standing probe
of positron production and transport in the Milky Way. Its line profile,
positronium continuum, and spatial distribution constrain the thermal and
ionization conditions in which positrons slow and annihilate
\citep{Jean2006,Prantzos2011,Siegert2016}. Wide-field measurements have
established a bright Galactic bulge and weaker disk, while a balloon-borne
Compton telescope has independently detected the bulge emission
\citep{Kierans2020}. These observations do not by themselves determine whether
part of the emission is concentrated in unresolved sources or compact central
structure. A search with INTEGRAL/IBIS found no compact Galactic line source and
reported a $2\sigma$ limit near $1.6\times10^{-4}\phflux$ for a roughly 10\,Ms
Galactic-centre exposure \citep{DeCesare2011}. A pointed instrument reaching
well below this scale would provide a complementary test of localized
annihilation emission.

Laue diffraction optics offer a way to concentrate a narrow gamma-ray band
onto a small focal plane \citep{Barriere2009}. Focusing reduces the
detector area exposed to backgrounds for a given collecting area, but its
benefit depends on the background generated within and immediately around the
focal plane \citep{Weidenspointner2006}. Stacked transition-edge-sensor (TES)
microcalorimeters are attractive at 511\keV because their fine energy response
supports a narrow line window and line-profile measurements
\citep{Gottardi2021,Shirazi2023,Hu2026}. The same dense absorbers and nearby
cryogenic structures that provide stopping power and thermal control can,
however, create prompt energy deposits and radioactive daughters. Passive mass
added close to the detector may therefore exchange one background mechanism
for another.

This work tests a direct geometric remedy. The TES focal plane is moved
laterally away from the projection of the central cold assembly and placed in a
side chimney with a compact BGO anticoincidence enclosure. The question is not
whether one shield material attenuates a selected particle in isolation. It is
whether the complete detector--cryostat arrangement lowers prompt and delayed
counts in the science window without sacrificing focused photons. We answer
that question by comparing two complete assemblies under identical source,
response, selection, and mission definitions. The design comparison is based on
existing transport and activation data; no component-wise transport ablation is
inferred.

Section~\ref{sec:instrument} describes the optics, TES array, and compared
geometries. Section~\ref{sec:methods} defines the atmospheric fields,
activation normalization, event response, selection, and mission fold.
Section~\ref{sec:results} presents the physical origin of the background
change, the matched cut flow, and the resulting sensitivity. The implications
and limits of the comparison are discussed in Section~\ref{sec:discussion}.

\section{Instrument and compared geometries}
\label{sec:instrument}

\subsection{Focused signal and TES focal plane}

The instrument concept combines a 10\,m-focal-length Ge(111) Laue optic with a
stacked tantalum-absorber TES focal plane. The ray set used here began with
150,000 photons incident on an illuminated crystal aperture of
20.08476\,cm$^2$. Diffraction and downstream geometric acceptance produced
37,194 focal-plane event-list entries. The same accepted-ray population was
transported through both detector assemblies. This construction isolates the
effect of the focal-plane environment from the upstream optical calculation.
The optics treatment follows established Laue-lens concepts, while particle
transport and detector interactions are represented with Geant4 and MEGAlib
\citep{Barriere2009,Agostinelli2003,Allison2016,Zoglauer2006}.

The absorber array contains six layers, each with 376 tantalum pixels, for a
total of 2256 pixels. Each pixel is
$1.5\times1.5\times3.0$\,mm$^3$, and adjacent layers are separated by 12\,mm.
These dimensions and the optical input are fixed in the comparison. The
reported effective-area change is therefore produced by the surrounding
assembly and the shared event selection, not by changing the TES inventory or
the incident ray sample.

\subsection{Reference and chimney geometries}
\label{sec:geometries}

In the reference assembly, the TES stack lies beneath the projected footprint
of the central cold stages and their support structures. This arrangement
provides a compact axial package, but it places activated cryogenic material
close to the line-sensitive absorber. The reference retains its established
active anticoincidence enclosure and serves as an end-to-end performance
baseline.

In the optimized assembly, the entire TES stack is translated laterally by
35.55\,cm and raised by 2.40\,cm into a side chimney. The chimney has aluminum
inner and outer radii of 10.75 and 11.05\,cm, corresponding to a 3\,mm wall. Its
lower end is separated from the adjacent cold-stage boundary by 2.5\,mm. A
6.0\,cm upper optical opening preserves the focused beam, and 1.5\,cm cold-stage
ports accommodate required penetrations. A tungsten frame surrounding the
beam has 6.0 and 5.4\,cm outer and inner square dimensions and a 2.0\,cm axial
depth.

Three active BGO volumes enclose the chimney focal plane. The cylindrical side
shield is 40\,mm thick radially and 130\,mm long. Front and rear BGO end pieces
are each 40\,mm thick; the front piece contains a square optical recess, and the
rear piece contains a circular cold-service opening. The transport geometry
uses an 80\keV native detector trigger. Event analysis applies the more inclusive
criterion that the summed BGO energy associated with the TES event be strictly
below 50\keV. The lower offline threshold is used uniformly for the reported
chimney results.

Figure~\ref{fig:geometry} summarizes the geometric mechanism tested here. It is
a sectional schematic rather than a fabrication drawing. Its relevant feature
is the separation of the active TES volume from the central cold-stage
projection while preserving a continuous BGO enclosure around the new focal
plane location.

\begin{figure}[t]
  \centering
  \IfFileExists{figures/fig_geometry_causal_section.pdf}{%
    \includegraphics[width=0.98\textwidth]{figures/fig_geometry_causal_section.pdf}%
  }{%
    \fbox{\parbox[c][5.5cm][c]{0.93\textwidth}{\centering
    Publication figure pending: reference and chimney geometry with selected
    delayed-event origin fractions.}}%
  }
  \caption{Sectional comparison of the reference assembly and the
  chimney-mounted TES assembly. The plotted percentages are the fractions of
  selected delayed events whose production positions lie in the central
  cryogenic-stage region. The physical comparison is between complete
  detector assemblies.}
  \label{fig:geometry}
\end{figure}

\section{Simulation and analysis}
\label{sec:methods}

\subsection{Atmospheric prompt fields}

The prompt radiation field is based on the EXPACS/PARMA atmospheric model
\citep{Sato2015,Sato2016}. The reference environment is latitude
$34^{\circ}$, longitude $100^{\circ}$, altitude 38\,km, vertical cutoff
rigidity 11.6\,GV, solar modulation index 118.3, and date 2025-08-31. Eight
families are transported: alpha particles, electrons, positrons, photons,
negative and positive muons, neutrons, and protons. Each family is divided into
20 equal bins in the cosine of zenith angle over the full sphere. Particles are
generated on a 60\,cm-radius far-field surface.

The full-sphere integral fluxes are 4.79966 for photons, 0.462239 for neutrons,
0.199002 for electrons, 0.116981 for positrons, 0.112301 for protons, 0.0114871
for alpha particles, 0.00496893 for negative muons, and 0.00557001 for positive
muons, all in cm$^{-2}$\,s$^{-1}$. The spectral axes supplied to transport are
total kinetic energy in keV. Thus a tabulated non-alpha energy in MeV is
multiplied by 1000, whereas an alpha-particle energy in MeV per nucleon is
multiplied by $4\times1000$. The 160 angular spectra are individually
normalized. The photon field is the total broadband EXPACS/PARMA component,
including its tabulated annihilation bump.

For prompt family $f$, accepted events carry a common normalization

\begin{equation}
  w_f^{\mathrm{prompt}}=\left(\sum_j T_{fj}\right)^{-1},
  \label{eq:promptweight}
\end{equation}

where $T_{fj}$ is the simulated live-time equivalent of angular job $j$.
Rates are sums of these event weights after each selection. The chimney
transport sample contains 114.99 million primaries across the eight families,
including 112.42 million photons. The reference sample has lower prompt
statistics, but all rate comparisons use the same physical normalization.

\subsection{Activation and delayed events}

Prompt irradiation also produces radioactive nuclei in detector and cryostat
materials. The inventory is evolved to a 15-day mature state using the time
profile of each incident family. Nuclear identities and ground-state handling
follow NUBASE2020 \citep{Kondev2021}. This bottom-up treatment follows the
general strategy used for balloon gamma-ray-instrument background studies
\citep{Gallego2025}. Production records retain the incident
family, parent nuclide, material, and position. Delayed decays are sampled at
their production positions, allowing selected detector events to be traced
back to the structure in which the parent was created. The maximum accepted
position-matching distance is $10^{-5}$\,cm.

One million decays are transported for each family, for eight million delayed
histories per geometry. A selected delayed event from family $f$ has weight

\begin{equation}
  w_f^{\mathrm{delayed}}(t)=\frac{A_f(t)}{N_f^{\mathrm{decay}}},
  \label{eq:delayedweight}
\end{equation}

where $A_f(t)$ is the activity produced by that family and
$N_f^{\mathrm{decay}}=10^6$. At day 15, the summed activities are 1352.30\,Bq
for the reference assembly and 211.387\,Bq for the chimney assembly. Activity
alone is not treated as a science-window background predictor: the parent
location, decay emission, energy response, veto, and topology selection all
remain necessary.

\subsection{Detector response and event selection}
\label{sec:selection}

We assume independent Gaussian energy smearing for each TES pixel with a full
width at half maximum of 0.420\keV, corresponding to
$\sigma=0.17836\keV$. Hits below 0.300\keV after smearing are discarded. All
primary results use the inclusive line window 510.58--511.42\keV. TES and BGO
deposits within 1\,$\mu$s are associated with the same event. A chimney event
passes anticoincidence only if the summed associated BGO energy is below
50\keV.

The final topology test uses the optical entrance as a directional constraint.
A single-site TES event is retained because no Compton direction can be
reconstructed. For two-site events, both interaction orders are tested, and an
event is accepted if at least one physical Compton cone intersects the entrance
aperture. For three to six sites, all admissible interaction sequences are
enumerated and scored by their Compton-kinematic residuals; equal residuals are
resolved in favor of the longer first lever arm, and the best sequence must
intersect the aperture. Events with more than six sites are retained to
avoid imposing an unvalidated high-multiplicity rejection. Nine representative
positions per pixel, comprising its eight corners and centroid, give 81
position pairs for the first two interactions. Each reconstructed cone is
sampled at 24 azimuths for the aperture-intersection calculation.

The focused-signal effective area after selection is

\begin{equation}
  \aeff=A_{\mathrm{ap}}\frac{N_{\mathrm{sel}}}{N_{\mathrm{in}}},
  \label{eq:aeff}
\end{equation}

where $A_{\mathrm{ap}}=20.08476$\,cm$^2$ and
$N_{\mathrm{in}}=37{,}194$. Its statistical uncertainty is evaluated as a
binomial counting uncertainty propagated through Eq.~\eqref{eq:aeff}.

\subsection{Trajectory fold and sensitivity}
\label{sec:mission}

Mission performance is evaluated on a synthetic 20-day balloon trajectory with
81 nodes separated by 0.25\,d. The source remains at $45^{\circ}$ elevation.
The trajectory spans an altitude envelope of $38.75\pm2.50$\,km, latitude
$34.00^{\circ}\pm0.25^{\circ}$, and longitude
$100.00^{\circ}\pm0.25^{\circ}$. Atmospheric depth and geomagnetic cutoff vary
consistently with those coordinates. Five transport anchors at
days 0, 5, 10, 15, and 20 set the time-dependent detector rates; interpolation
between anchors provides the node rates. The reference anchors contain 20\,ks
of replay exposure each, whereas the higher-statistics chimney anchors contain
200\,ks each. Their physical definitions and selections are matched, and the
finite replay uncertainty is propagated separately. The focused signal is attenuated by
the atmospheric slant transmission and by accidental anticoincidence. The
511\keV transmission is 0.65203 at days 15 and 20. It is calculated in a
plane-parallel dry-air model as
$\tau_{511}=\exp[-\mu_{\mathrm{air}}X_{\mathrm{vert}}/\cos45^{\circ}]$, with
$\mu_{\mathrm{air}}=0.08736$\,cm$^2$\,g$^{-1}$ from XCOM data near 511\keV
\citep{NISTXCOM}. At day 15 the simulated
accidental-veto survival is 0.995341 in the chimney assembly and 0.967071 in the
reference. At each anchor, prompt and delayed streams are represented as
independent Poisson arrivals, adjacent events within 1\,$\mu$s are merged, and
the response and selections are reapplied. Background-rate ratios and signal
accidental-survival factors are linearly interpolated between anchors; the
cumulative counts are integrated with the trapezoidal rule.

For a top-of-atmosphere line flux $F$, the cumulative signal and background are

\begin{align}
  N_s(F,T) &= F\int_0^T \aeff(t)\,\tau_{511}(t)\,
                   \epsilon_{\mathrm{acc}}(t)\,dt, \\
  N_b(T) &= \int_0^T R_b(t)\,dt,
  \label{eq:counts}
\end{align}

where $\tau_{511}$ is atmospheric transmission and
$\epsilon_{\mathrm{acc}}$ is the accidental-veto survival. Two significance
conventions are reported. The Gaussian counting approximation is
$Z_G=N_s/\sqrt{N_b}$, and the background-known Asimov expression is

\begin{equation}
  Z_A=\left\{2\left[(N_s+N_b)\ln\left(1+\frac{N_s}{N_b}\right)-N_s\right]
       \right\}^{1/2}
  \label{eq:asimov}
\end{equation}

\citep{Cowan2011}. Equation~\eqref{eq:asimov} is the known-background,
single-bin Asimov median discovery significance. The line-flux threshold at a
specified $Z$ is obtained by solving the corresponding expression for $F$.
Neither expression profiles nuisance parameters or includes a
look-elsewhere correction. The threshold is a counting-statistics metric, not
a confidence interval that includes instrumental or atmospheric systematics.

Independent uncertainties from weighted transport counts, trajectory replay,
effective area, and accidental-veto probes are propagated with the delta
method. For the chimney assembly their relative one-standard-deviation
contributions are 18.67\% for background transport, 1.13\% for trajectory
replay, 0.30\% for effective area, and 0.002\% for the accidental probe. The
combined background uncertainty is 18.71\%; because the weak-signal Gaussian
threshold scales approximately as $\sqrt{N_b}$, the corresponding relative
uncertainty of the 3$\sigma$ flux threshold is 9.36\%.

\subsection{Use of generative artificial intelligence}

An OpenAI language model assisted with English-language editing and manuscript
assembly. It did not generate transport or activation data. The authors retain
responsibility for the numerical values, scientific interpretation, and final
text.

\section{Results}
\label{sec:results}

\subsection{Geometric decoupling of delayed background}
\label{sec:coupling}

The clearest causal diagnostic is the production position of delayed events
that survive all science selections. In the reference assembly, 32.18\% of the
selected delayed rate originates in the central detector-refrigeration and
cold-stage region. In the chimney assembly this fraction is 3.92\%. The
corresponding rates are 0.01214 and $1.42\times10^{-4}\cps$, respectively. This
factor of about 85 in rate is larger than the change in total activation, which
shows that the improvement is not explained by inventory reduction alone.
The result is consistent with the combined relocated-TES and BGO geometry
reducing the probability that a decay photon reaches the TES and survives the
active veto. Because several structural and active elements differ, this is not
a single-variable causal measurement.

The source distribution also changes qualitatively. TES-near structures account
for 47.29\% of the selected delayed rate in the reference and 89.38\% in the
chimney assembly. This rise in fractional share does not indicate a higher
TES-near rate; it indicates that the remote cold-stage contribution has largely
been removed from the residual. The largest identified volume-specific group
in the chimney geometry is $^{62}$Cu produced in a copper support near the TES,
contributing $4.22\times10^{-4}\cps$, or 11.6\% of the delayed total. The
remaining optimization target is therefore local support and readout material,
not the displaced central cold assembly.

\subsection{Matched event cut flow}
\label{sec:cutflow}

Table~\ref{tab:cutflow} gives the simulated cut flow for both assemblies. Before
anticoincidence, the chimney assembly has a prompt line-window rate of
2.235\cps and a delayed rate of 0.1275\cps. The 50\keV BGO condition reduces
these values to 0.005794 and 0.004019\cps. The topology condition removes no
additional prompt events in the finite final sample and reduces the delayed
rate to 0.003625\cps, leaving a total of 0.009419\cps. Thus the BGO condition is
responsible for nearly all prompt rejection and most delayed rejection, whereas
the directional topology test supplies a smaller, physically distinct cleanup
of multi-site delayed events.

\begin{table}[t]
  \centering
  \caption{Matched fixed-source-state line-window cut flow. Rates are direct
  transport estimates before trajectory interpolation. Each column uses the
  full active-veto logic of its complete assembly; Figure~\ref{fig:cutflow}
  separately isolates the active-BGO stage. Signal effective areas use the same
  37,194 focused input photons.}
  \label{tab:cutflow}
  \small
  \begin{tabularx}{\textwidth}{@{}Yrrrr@{}}
    \toprule
    & \multicolumn{2}{c}{Reference assembly}
    & \multicolumn{2}{c}{Chimney assembly} \\
    \cmidrule(lr){2-3}\cmidrule(lr){4-5}
    Quantity & Selected records & Rate or area & Selected records & Rate or area \\
    \midrule
    Prompt, before active veto & 1083 & 5.79018\cps & 3212 & 2.23515\cps \\
    Prompt, after assembly veto & 7 & 0.018704\cps & 12 & 0.005794\cps \\
    Prompt, final topology & 6 & 0.016032\cps & 12 & 0.005794\cps \\
    Delayed, before active veto & 1295 & 0.169494\cps & 3249 & 0.127481\cps \\
    Delayed, after assembly veto & 436 & 0.043412\cps & 122 & 0.004019\cps \\
    Delayed, final topology & 394 & 0.038266\cps & 111 & 0.003625\cps \\
    Total background, final & 400 & 0.054298\cps & 123 & 0.009419\cps \\
    Focused signal, final & 28,006 & 15.12324\,cm$^2$ & 27,936 & 15.08544\,cm$^2$ \\
    \bottomrule
  \end{tabularx}
\end{table}

The reference and chimney rates cannot be read as a single-shield attenuation
measurement because the complete assemblies differ. They do establish the
system-level endpoint under a shared source and analysis. The direct final
background falls by 82.7\%, while the focused effective area falls by only
0.25\%. The ratio of final effective areas is 0.99750. Binomial counting gives
15.123$\pm$0.045 and 15.085$\pm$0.045\,cm$^2$ for the reference and chimney
assemblies. The difference is therefore small in both absolute and mission
terms.

Figure~\ref{fig:cutflow} displays the rate change on a logarithmic scale. In the
chimney result, the final prompt component is entirely represented by 12
weighted photon events. The final delayed component contains 111 events spread
over several incident-family weights. The low effective number of weighted
survivors, especially for prompt photons, is why the main sensitivity result
retains a transport-counting uncertainty despite the large number of incident
histories. The reference prompt endpoint contains only six raw photon survivors
and has a 40.8\% relative component uncertainty.

\begin{figure}[t]
  \centering
  \IfFileExists{figures/fig_matched_cutflow_performance.pdf}{%
    \includegraphics[width=0.94\textwidth]{figures/fig_matched_cutflow_performance.pdf}%
  }{%
    \fbox{\parbox[c][5.5cm][c]{0.88\textwidth}{\centering
    Publication figure pending: matched prompt and delayed cut-flow rates.}}%
  }
  \caption{Total prompt-plus-delayed 510.58--511.42\keV rates before
  anticoincidence, after the active BGO veto, and after final topology
  selection. Points show central weighted estimates; connecting lines guide
  the eye. The different active enclosures are properties of the two complete
  assemblies.}
  \label{fig:cutflow}
\end{figure}

\subsection{Trajectory background and line sensitivity}
\label{sec:sensitivity}

Trajectory folding preserves the geometry advantage. At day 15 the mature
background rate is 0.051280\cps for the reference and 0.009280\cps for the
chimney assembly, an 81.9\% reduction. For a top-of-atmosphere line flux
$F_0=10^{-4}\phflux$, the corresponding signal rates are
$9.536\times10^{-4}$ and $9.790\times10^{-4}\cps$. The slightly higher
chimney signal rate, despite its marginally smaller effective area, results from
the higher simulated accidental-veto survival at that trajectory state.

Over 20 days, the reference accumulates 86,806 background counts and 1635
signal counts at $F_0$. The chimney assembly accumulates 15,552 background
counts and 1678 signal counts. These counts give Gaussian significances of 5.55
and 13.46 and Asimov significances of 5.53 and 13.22. The chimney assembly
reaches 3$\sigma$ after 0.724\,d in the Gaussian approximation and 0.753\,d with
the Asimov expression. It reaches 5$\sigma$ after 1.94 and 2.03\,d,
respectively. The reference requires 5.52\,d for 3$\sigma$ and nearly 16\,d for
5$\sigma$.

\begin{table}[t]
  \centering
  \caption{Twenty-day performance for the two complete detector assemblies.
  Signal quantities use $F_0=10^{-4}\phflux$. Quoted threshold uncertainties
  contain transport, trajectory-replay, effective-area, and accidental-probe
  counting terms only.}
  \label{tab:mission}
  \small
  \begin{tabularx}{\textwidth}{@{}Yrr@{}}
    \toprule
    Metric & Reference assembly & Chimney assembly \\
    \midrule
    Day-15 background rate (s$^{-1}$) & 0.051280 & 0.009280 \\
    Day-15 signal rate (s$^{-1}$) & $9.536\times10^{-4}$ & $9.790\times10^{-4}$ \\
    Cumulative background counts & 86,806 & 15,552 \\
    Cumulative signal counts & 1635 & 1678 \\
    Gaussian significance & 5.55 & 13.46 \\
    Asimov significance & 5.53 & 13.22 \\
    Gaussian time to 3$\sigma$ (d) & 5.52 & 0.724 \\
    Asimov time to 3$\sigma$ (d) & 5.55 & 0.753 \\
    Gaussian time to 5$\sigma$ (d) & 15.93 & 1.94 \\
    Asimov time to 5$\sigma$ (d) & 16.01 & 2.03 \\
    Gaussian 3$\sigma$ flux ($\phflux$) &
      $(5.406\pm0.401)\times10^{-5}$ &
      $(2.229\pm0.209)\times10^{-5}$ \\
    Asimov 3$\sigma$ flux ($\phflux$) &
      $5.415\times10^{-5}$ & $2.238\times10^{-5}$ \\
    Gaussian 5$\sigma$ flux ($\phflux$) &
      $9.010\times10^{-5}$ & $3.716\times10^{-5}$ \\
    Asimov 5$\sigma$ flux ($\phflux$) &
      $9.035\times10^{-5}$ & $3.740\times10^{-5}$ \\
    \bottomrule
  \end{tabularx}
\end{table}

The Gaussian 3$\sigma$ flux threshold improves from
$(5.406\pm0.401)\times10^{-5}$ to
$(2.229\pm0.209)\times10^{-5}\phflux$. The ratio is 0.412, corresponding to a
factor of 2.43 improvement. Asimov thresholds differ from the Gaussian values
by less than 0.5\% at 20 days, as expected when the accumulated background is
large. At shorter exposure the Asimov curve is modestly more conservative.
Figure~\ref{fig:mission} shows the cumulative significance and flux threshold.

\begin{figure}[t]
  \centering
  \IfFileExists{figures/fig_mission_significance_sensitivity.pdf}{%
    \includegraphics[width=0.98\textwidth]{figures/fig_mission_significance_sensitivity.pdf}%
  }{%
    \fbox{\parbox[c][5.8cm][c]{0.92\textwidth}{\centering
    Publication figure pending: cumulative significance and line-flux threshold
    flux over the synthetic trajectory.}}%
  }
  \caption{Mission-folded line performance for a continuously observed source
  at $45^{\circ}$ elevation. Left: cumulative Asimov significance for
  $F_0=10^{-4}\phflux$. Right: Asimov 3$\sigma$ line-flux threshold. The
  trajectory is a synthetic instrument benchmark, not a flight visibility
  schedule. The grey IBIS scale is a linear Gaussian conversion of the
  published $2\sigma$ limit and is not a reported IBIS $3\sigma$ result.}
  \label{fig:mission}
\end{figure}

\section{Discussion}
\label{sec:discussion}

\subsection{Why the chimney geometry works}

The result is consistent with two linked mechanisms. First, the lateral
translation weakens
the solid-angle coupling between the TES and activated material in the central
cold stages. The 32.18\% to 3.92\% change in the selected delayed-source
fraction provides a direct event-lineage measure of that decoupling. Second,
the BGO enclosure is centered on the translated TES stack. It rejects events
that reach the focal volume with accompanying energy deposition while leaving
the optical opening unobstructed. The result is not simply proportional to
total activity: day-15 activity decreases by a factor of 6.4, whereas the
selected central-cold-stage delayed rate decreases by about 85. Geometry and
event topology determine how much of the inventory becomes line-window
background.

The signal result is equally important. A shield configuration that lowers
background by blocking the focused beam would not improve a pointing
telescope. Here the final effective-area ratio is 0.9975, and the 20-day signal
count is slightly higher after accidental-veto folding. The line sensitivity
gain is therefore controlled by background suppression rather than a trade of
signal acceptance for rejection. In the weak-source Gaussian limit,
$F_{3\sigma}\propto\sqrt{N_b}/K_s$, where $K_s$ is the signal-count kernel. The
simulated background and signal-kernel ratios predict the matched factor of
2.43 to rounding.

The residual source map also gives a disciplined next step. Once the central
cold-stage contribution is geometrically suppressed, most selected delayed
background is associated with material close to the TES. Copper activation,
including $^{62}$Cu, is a measurable contributor. Future mechanical refinement
should therefore examine the minimum necessary copper support mass, its
placement relative to the active BGO, and material substitutions compatible
with cryogenic and structural requirements. Those changes require dedicated
transport and activation comparisons; the present data do not quantify their
benefit.

\subsection{Implications for a pointed 511 keV telescope}

The 20-day chimney threshold of $2.23\times10^{-5}\phflux$ is about seven times
below the published $2\sigma$ IBIS compact-source limit of
$1.6\times10^{-4}\phflux$ \citep{DeCesare2011}, although the instruments,
confidence conventions, exposure histories, and fields of view are different.
The comparison is therefore a scale indicator rather than a forecast of a
specific source detection. Within the present idealized mission fold, a source
at $10^{-4}\phflux$ reaches the quoted significance rapidly under the
continuous-exposure, known-background assumptions, leaving statistical margin
for subdivision by time or energy. The assumed TES response of 0.420\keV FWHM
would also support line-centroid and profile studies that are inaccessible to a
broad window alone \citep{Jean2006,Siegert2016,Hu2026}.

The reference-to-chimney ratio is less sensitive to shared modeling choices
than either absolute sensitivity value. The atmospheric
families, total-energy axes, focused rays, response function, science window,
directional selection, and trajectory are matched. Their ratio consequently
isolates the detector-assembly change more cleanly than an absolute rate can.
The comparison nevertheless remains between complete geometries; it cannot
assign a unique fraction of the improvement to lateral displacement, BGO
thickness, or aperture details.

\subsection{Limitations}
\label{sec:limitations}

The modeled domain ends at the detector and cryostat. It does not include a
complete Laue-lens structure, balloon gondola, pressure vessel, flight train,
telemetry hardware, or their prompt and activation backgrounds. The synthetic
trajectory assumes continuous on-source exposure at fixed elevation and does
not impose Earth occultation, target visibility, attitude excursions,
dead-time losses, detector downtime, or a realistic observing schedule.
Absolute mission sensitivity will change when these elements are added.

The BGO response is an ideal deposited-energy veto. Light collection,
photodetector response, channel thresholds, timing tails, pile-up, dead time,
and threshold nonuniformity are not modeled. The 50\keV offline condition is
more inclusive than the 80\keV native transport trigger, so hardware realization
must demonstrate the lower analysis threshold with acceptable noise and timing.
Likewise, the TES response is represented by an independent Gaussian width and
a sharp post-noise threshold. Calibration tails, common-mode noise,
nonlinearity, pixel-to-pixel resolution variation, and event reconstruction
failures are outside the present uncertainty.

The reported errors are statistical. They include finite weighted transport,
trajectory replay, focused-ray acceptance, and accidental-veto probe counts,
but not atmospheric-model uncertainty, nuclear production cross sections,
decay data, material composition, geometry tolerances, or detector-response
systematics. Weighted prompt survivors are sparse even after the expanded
photon sample; exact component intervals are correspondingly asymmetric and
can be much wider than the central rate. The quoted 9.36\% error on the chimney
3$\sigma$ threshold should not be interpreted as total mission uncertainty.

Finally, the chimney layout has passed geometric construction and overlap
checks, but it is not mechanically qualified. Thermal conductance, vibration,
weld design, BGO segmentation and support, cable routing, magnetic shielding,
and cryogenic service loads require engineering closure. The performance gain
reported here motivates that work; it does not replace it.

\section{Conclusions}
\label{sec:conclusions}

Moving the TES focal plane into a lateral chimney and enclosing it with active
BGO suppresses the modeled day-15 detector--cryostat background by 81.9\%
while retaining 99.75\% of the focused effective area. The matched 20-day fold
improves the Gaussian 3$\sigma$ 511\keV line-flux threshold by a factor
of 2.43, to $(2.229\pm0.209)\times10^{-5}\phflux$.

The production-position record identifies the physical reason: selected
delayed events from the central cold assembly fall from 32.18\% to 3.92\%, so
the residual becomes concentrated in TES-near structures that can be targeted
by later material and mechanical studies. These results support the chimney
TES and BGO enclosure as the detector-geometry baseline for a more complete
payload simulation and engineering qualification.

\section*{Data availability}

The aggregate numerical inputs underlying the tables and figures are included
with the manuscript source package. Event-level transport and activation
products are available from the authors upon reasonable request. Publicly
documented physics packages and data evaluations are cited in the text.

\section*{Statements and declarations}

\noindent\textbf{Competing interests.} The authors declare no competing
interests.

\noindent\textbf{Ethics approval and consent to participate.} Not applicable.

\noindent\textbf{Funding and author contributions.} Details are withheld in
this review manuscript and will be supplied with author identities.

\bibliographystyle{plainnat}
\bibliography{references}

\end{document}
